\documentclass[11pt]{amsart}

\usepackage[margin=1in]{geometry}
\usepackage{amsmath,amssymb,amsthm,mathtools}
\usepackage{enumitem}
\usepackage{xcolor}

\newtheorem{theorem}{Theorem}
\newtheorem{proposition}[theorem]{Proposition}
\newtheorem{lemma}[theorem]{Lemma}
\newtheorem{remark}[theorem]{Remark}
\newtheorem{conjecture}[theorem]{Conjecture}

\newcommand{\E}{\mathbb E}
\newcommand{\N}{\mathbb N}
\newcommand{\R}{\mathbb R}
\newcommand{\Span}{\operatorname{span}}
\newcommand{\supp}{\operatorname{supp}}

\title{A characterization of independent random variables which perform stable phase retrieval}
\author{Pedro Abdalla\textsuperscript{1}}

\author{Jaume de Dios Pont\textsuperscript{2}}

\author{Jo\~ao Pedro Ramos\textsuperscript{3}}

\author{Mitchell A.~Taylor\textsuperscript{4}}

\begin{document}
\begin{abstract}
We provide a complete characterization of when a subspace of $L^2(\mu;\mathbb{R})$ spanned by independent mean-zero
random variables performs stable phase retrieval. More specifically, we prove that stable phase retrieval holds for such a subspace if all but one random variable satisfy a uniform two-sided $L^1$ bound. This confirms a conjecture by Calderbank,
Daubechies, Freeman, and Freeman. %, who were able to characterize when the span of independent mean-zero
%random variables, each summand with a disjoint indicator function, performs stable phase retrieval in $L^2(\mathbb{R};\mathbb{R})$.

Our
argument is based on a quantitative compactness scheme: after applying an orthogonal reduction, we construct 
a counterexample sequence which is decomposed into a convergent head component and a
diffuse tail component. The head is recovered from the quadratic observable
$|X|^2-|Y|^2$ by bounded probe functions, while the tail is controlled by invoking the weak-limit of triangular arrays. %through
%weak-limit rigidity and a uniform lower $L^1$-$L^2$ estimate for independent
%sums.
\end{abstract}
\maketitle
\markboth{ABDALLA, DE DIOS PONT, RAMOS, AND TAYLOR}{STABLE PHASE RETRIEVAL FOR INDEPENDENT RANDOM VARIABLES}
\begingroup
\renewcommand{\thefootnote}{\arabic{footnote}}
\footnotetext[1]{Department of Mathematics, University of California, Irvine, California, USA}
\footnotetext[2]{Center for Data Science, New York University, New York, New York, USA}
\footnotetext[3]{Instituto de Matem\'atica Pura e Aplicada (IMPA), Rio de Janeiro, RJ, Brazil}
\footnotetext[4]{Department of Mathematics, ETH Z\"urich, Z\"urich, Switzerland}
\endgroup
\setcounter{footnote}{0}

\section{Introduction}

The phase retrieval problem is concerned with recovering a vector or function
from phaseless measurements. In the real-valued setting, this means recovery up to a multiplicative
sign. Although the finite-dimensional theory is by now extensive; see, for instance,
\cite{bandeira2014saving,eldar2014stability} and the survey
\cite{grohs2020survey}, the infinite-dimensional case is much more
delicate. For example, phase retrieval itself can hold in infinite-dimensional Hilbert
spaces \cite{cahill2016infinite}, but uniform stability fails in the general
setting of continuous frames \cite{alaifari2017continuous,alaifari2025cheeger}, and natural
measurement models may be ill-posed \cite{alaifari2021gabor}. On the other hand, there are non-trivial
structured subsets of reproducing kernel Hilbert spaces that perform phase retrieval with H\"older stability bounds \cite{cahill2016infinite,freeman2025cahill}.

We adopt the function-space point of view from \cite{FOTP}. Let $(X,\|\cdot\|_X)$ be a real
normed function space and $E \subset X$ be a subspace, we say that $E$ does stable
phase retrieval if there exists a constant $C < \infty$ such that
\[
\min_{\sigma \in \{\pm 1\}} \|f-\sigma g\|_X
\leq C \, \bigl\||f|-|g|\bigr\|_X
\qquad
\text{for all } f,g \in E.
\]
It is natural to ask which subspaces of a given function
space can do stable phase retrieval. The main results of \cite{christ2022examples, FOTP,garcia2025isometric,garcia2025existence} show that
many Banach lattices admit infinite-dimensional closed subspaces
with this property, and, importantly, that such subspaces share remarkable structural properties. In particular, the authors in \cite{FOTP} prove an orthogonal reduction
in Hilbert spaces: to verify stable phase retrieval on a subspace, it is
enough to only check it for orthogonal pairs $(f,g)$. Moreover, for general Banach lattices, it suffices to prove stable phase retrieval for ``well-separated" vectors. This reduction is the starting point of the present
 argument.

In \cite{calderbank2022stable}, Calderbank, Daubechies, Freeman, and Freeman
proved the following influential theorem.

\begin{theorem}\cite[Theorem Number?]{calderbank2022stable}
\label{ThmCDFF}
Let $(X_n)_{n=1}^\infty$ be an orthonormal sequence of independent mean-zero
random variables in $L^2([0,1])$. The following statements are equivalent.
\begin{enumerate}[label=\rm(\arabic*)]
\item The closed span of $(X_n+1_{(n,n+1]})_{n=1}^\infty$ does stable phase
retrieval in $L^2(\R)$.
\item The $L^1$ and $L^2$ norms of $(X_n)$ are uniformly comparable; that is,
there exists $A>0$ such that
\[
A\|X_n\|_{L^2([0,1])}\leq \|X_n\|_{L^1([0,1])}\leq \|X_n\|_{L^2([0,1])}
\qquad
\text{for all } n\in \N.
\]
\end{enumerate}
\end{theorem}

Following that work, the authors communicated the following conjecture. \textcolor{red}{any formal reference for the conjecture?}

\begin{conjecture}\label{conj:cdff}
If $(X_n)$ is an orthonormal sequence of independent mean-zero random
variables in $L^2$, then $\overline{\Span}(X_n)$ does stable phase retrieval in $L^2$ if
and only if there exist constants $0<A\leq B<1$ such that
\[
A\leq \|X_n\|_{L^1}\leq B
\]
for all but at most one $n\in \N$.
\end{conjecture}

The necessity of the conditions in Conjecture \ref{conj:cdff} is well
known in the literature. Indeed,
by \cite{FOTP}, if a subspace of $L^p(\mu)$
does stable phase retrieval, then the $L^1$ and $L^p$ norms are equivalent on
that subspace, so the lower bound $A>0$ is necessary. The upper bound $B<1$
reflects the constraint that a centered random variable with
$\|f\|_{L^2}=1$ satisfies
$\|f\|_{L^1}=\|f\|_{L^2}$ if and only if it is Rademacher-valued (random sign), and the
closed span of the Rademacher functions clearly violates phase retrieval.
Since at most one coordinate can be Rademacher, a uniform gap $B<1$ is needed
for all remaining coordinates. The main result of this note shows that such conditions are also sufficient to establish stable phase retrieval.


There is an interesting connection between Theorem~\ref{ThmCDFF} and the Pauli phase retrieval problem.
Let $G$ be a locally compact Abelian group. We say that a subspace $E\subseteq L^2(G)$ does \emph{Pauli phase retrieval}
if for all $f,g\in E$ we have $|f|=|g|$ and $|\mathcal{F}f|=|\mathcal{F}g|$ if and only if $f=\lambda g$ for some unimodular
scalar $\lambda$. The original \emph{Pauli problem} asked if, when $G=\mathbb{R}$, the whole space $E=L^2(\mathbb{R})$ does Pauli phase retrieval.
Since we now know that this is not the case, it has become an interesting problem to identify ``large" subspaces $E$ of $L^2(G)$ which perform Pauli phase 
retrieval and to prove stability estimates of the form
\[
\min_{|\lambda|=1} \|f-\lambda g\|_{L^2(G)}
\leq C \, \left(\bigl\||f|-|g|\bigr\|_{L^2(G)}+\bigl\||\mathcal{F}f|-|\mathcal{F}g|\bigr\|_{L^2(\widehat{G})}\right)
\qquad
\text{for all } f,g \in E.
\]
The connection between Theorem~\ref{ThmCDFF} and the Pauli phase retrieval problem can be seen as follows: Take $G=\mathbb{T}$ and $E$ to be the span of a sufficiently sparse subsequence $\Lambda$ of the characters $e^{2\pi i nx}$ (or, more conveniently, the real/imaginary parts $\sin(2\pi nx)$ or $\cos(2\pi nx)$  to avoid the obstruction $|e^{2\pi i nx}|=1$). In this case, if we expand $f=\sum_{n\in \Lambda} a_n \sin(2\pi nx)$ and $g=\sum_{n\in \Lambda} b_n \sin(2\pi nx)$ in the basis, then the absolute value of the Fourier transform allows us to obtain information on the magnitude of the coefficients $|a_n|$ and $|b_n|$, which is exactly what the  indicator functions in Theorem~\ref{ThmCDFF} had allowed us to do.  In other words, (with the caveat that we have replaced independent random variables with sparse Fourier series) the indicator functions in Theorem~\ref{ThmCDFF} give the same additional information as the Fourier transform does for the Pauli phase retrieval problem on $\mathbb{T}$. 

The above comparison can also be extended to the work \cite{christ2022examples}, where it was shown that if $(r_n)$ is an orthonormal sequence in $L^2(\mu)$ satisfying appropriate hypotheses (verified by sufficiently sparse sin series) then one can perform stable phase retrieval, i.e., one can stably reconstruct $f=\sum_n a_n r_n$ from $|f|$. If one were to add indicator functions to $r_n$ as in Theorem~\ref{ThmCDFF}, the question of stable phase retrieval for the resulting sequence $g_n=r_n+1_{(n,n+1]}$ would reduce to asking whether one can stably recover $f$ from $|f|$, together with knowledge of the sequence $(|a_n|)\in \ell^2.$ Evidently, this is a much easier task, and is one of the reasons why the proof of Theorem~\ref{conj:cdff} is much more complicated than that of Theorem~\ref{ThmCDFF}.




\textcolor{red}{Mention other nonlinearities. In particular, declipping \cite{abdalla2025sharp} and related references}.



\textcolor{blue}{Pedro: I suggest to reorder the next paragraph to appear before the Pauli PR}

Theorem \ref{thm:main} proves the sufficiency direction with
one exceptional coordinate. After relabeling the coordinates and normalizing
each variable in $L^2$, this is exactly the case in which all but at most one
variable satisfy uniform two-sided $L^1$ bounds. Our purpose here is not to
optimize constants, but to isolate a compactness mechanism behind our proof.
After applying the orthogonal reduction from \cite{FOTP}, we construct a counterexample sequence of pairs that splits
into a convergent head and a diffuse tail whose maximal coefficient tends to
zero. The head is identified from the relation $|X_n|^2-|Y_n|^2 \to 0$ by
means of bounded probes, while the tail is controlled through a weak-limit
rigidity argument and a uniform lower $L^1$-$L^2$ estimate for independent
sums. Together these ingredients force the two coefficient vectors to agree up
to sign, contradicting the orthogonal normalization. We now turn to the precise
formulation of the theorem and the contradiction setup. 

\textcolor{red}{TODO: put numbers for the sections}The rest of this manuscript is organized as follows. The next section
reduces the proof to a normalized counterexample array; the following two
sections develop the corresponding profile decomposition and auxiliary
estimates; and the final sections identify the limiting head and tail and
derive the contradiction.

\textcolor{red}{TODO: add notation section}
\section{Main result and contradiction setup}

We begin by stating the precise one-exception theorem proved in this note. We
then reformulate its failure in terms of a normalized orthogonal
counterexample sequence, which is the object analyzed in the rest of the
paper.

\begin{theorem}[Main Result]
\label{thm:main}
Let $\eta,\xi_2,\xi_3,\dots$ be independent real-valued random variables and
let $\delta > 0$ be such that
\[
\|\eta\|_{L^2} = 1,
\qquad
\E[\eta] = 0,
\]
and, for every $i \ge 2$,
\[
\|\xi_i\|_{L^2} = 1,
\qquad
\E[\xi_i] = 0,
\qquad
\delta \le \|\xi_i\|_{L^1} \le 1-\delta.
\]
Then there exists a constant $C_{\delta,\eta} \ge 1$ depending only on $\eta$ and $\delta$ such that, for every
$X,Y \in \overline{\Span}(\eta,\xi_i : i \ge 2)$,
\[
\min_{\sigma \in \{\pm 1\}} \|X-\sigma Y\|_{L^2}
\le
C_{\delta,\eta} \, \bigl\| |X|-|Y| \bigr\|_{L^2}.
\]
\end{theorem}

\begin{remark}
After relabeling the coordinates and normalizing each variable in $L^2$, this
is exactly the sufficiency direction of the formulation in which
\[
\delta \le \|\xi_i\|_{L^1} \le 1-\delta
\]
holds for all but at most one index.
\end{remark}

The rest of the proof is by contradiction. After the orthogonality
reduction from \cite{FOTP}, a failure of Theorem \ref{thm:main} produces a sequence of
orthogonal pairs with almost identical moduli. The remaining sections show
that no such sequence can exist.

\begin{proposition}
\label{prop:counterexample-array}
If the conclusion of Theorem \ref{thm:main} is false, then there exist:
\begin{enumerate}[label=\rm(\roman*)]
\item a triangular array $(\xi_{n,i})_{n \ge 1,\ i \ge 2}$ such that for each
fixed $n$ the family $(\eta,\xi_{n,i})_{i \ge 2}$ is independent and every
$\xi_{n,i}$ satisfies
\[
\|\xi_{n,i}\|_{L^2} = 1,
\qquad
\E[\xi_{n,i}] = 0,
\qquad
\delta \le \|\xi_{n,i}\|_{L^1} \le 1-\delta;
\]
\item vectors $a^{(n)}, b^{(n)} \in \ell^2(\N)$ such that
\[
\|a^{(n)}\|_2 = \|b^{(n)}\|_2 = 1,
\qquad
\langle a^{(n)}, b^{(n)} \rangle = 0,
\]
and for 
\[
X_n := a_1^{(n)}\eta + \sum_{i \ge 2} a_i^{(n)} \xi_{n,i},
\qquad
Y_n := b_1^{(n)}\eta + \sum_{i \ge 2} b_i^{(n)} \xi_{n,i},
\]
one has
\[
\bigl\| |X_n|-|Y_n| \bigr\|_{L^2} \longrightarrow 0.
\]
\end{enumerate}
\end{proposition}

\begin{proof}
If the uniform stable phase retrieval estimate fails, for any constant $C_n>0$ there exists a violating pair $(X_n,Y_n)$ that does not satisfy stable phase retrieval with a constant smaller or equal to $C_n$. By the orthogonal reduction
\cite[Theorem~3.9]{FOTP}, this pair can be replaced by an orthogonal pair
in the same two-dimensional span without increasing the quantity
$\bigl\| |X_n|-|Y_n| \bigr\|_{L^2}$. After normalizing both functions in the pair to have
$L^2$ norm one, the orthogonality between $X_n$ and $Y_n$ implies that
$\min_{\sigma=\pm 1}\|X_n-\sigma Y_n\|_{L^2}=\sqrt{2}$. Consequently, there exists a diverging sequence of positive numbers  $C_n$ for which
\begin{equation*}
\bigl\| |X_n|-|Y_n| \bigr\|_{L^2}< \frac{1}{C_n}\min_{\sigma=\pm 1}\|X_n-\sigma Y_n\|_{L^2} = \frac{\sqrt{2}}{C_n},
\end{equation*}
so in the limit as $n$ goes to infinity we must have $\bigl\| |X_n|-|Y_n| \bigr\|_{L^2}$ vanishes. After relabeling the good
variables row by row, we may write
\[
X_n = a_1^{(n)}\eta + \sum_{i \ge 2} a_i^{(n)} \xi_{n,i},
\qquad
Y_n = b_1^{(n)}\eta + \sum_{i \ge 2} b_i^{(n)} \xi_{n,i},
\]
where the exceptional coordinate $\eta$ is the same for every row. Because the random variables are centered, independent with unit $L^2$ norm, each
family $(\eta,\xi_i)_{i \ge 2}$ is orthonormal in $L^2$; hence each function $X_n$ is well-defined in $L^2$ and also
\begin{equation*}
    1= \|X_n\|_{L^2}^2 = \|a^{(n)}\|_2^2.
\end{equation*}
Similarly, we obtain that $\|b^{(n)}\|_2=1$. The orthogonality between $a^{(n)}$ and $b^{(n)}$ in $\ell_2$ follows from the orthogonality of $X_n$ and $Y_n$ in $L^2$. 
\end{proof}

\section{Profile decomposition}
Starting from the triangular array in Proposition \ref{prop:counterexample-array}, we reorder the vectors $a^{(n)}$ and $b^{(n)}$  so that the large coefficients in magnitude are mapped to the smaller indices. The core idea is to design a rearrangement such that part of the sequence $a^{(n)}$ converge strongly in $\ell_2$ and the rest of the sequence give rise to an infinitesimal triangular array. This 
head-tail decomposition governs the rest of the argument. The following lemma formalizes this decomposition.

\begin{lemma}[Reordered profile decomposition]\label{lem:profile}
Let $a^{(n)}$ and $b^{(n)}$ as in Proposition \ref{prop:counterexample-array}. There exists a permutation of the set $\{2,3,\dots\}$ applied to the coefficients and to the family
$(\xi_{n,i})_{i \ge 2}$ for which one may find 
a diverging sequence of positive integers $m_n$ such that, for
\[
a^{(n)}_{\mathrm{head}} := a^{(n)} \mathbf{1}_{\{1,\dots,m_n\}},
\quad
a^{(n)}_{\mathrm{tail}} := a^{(n)} \mathbf{1}_{\{m_n+1,m_n+2,\dots\}},
\]
and likewise for $b^{(n)}$, the following hold for a subsequence of $n$:
\begin{enumerate}
\renewcommand{\labelenumi}{\Roman{enumi}.}
\item There exists $a,b \in \ell^2(\mathbb{N})$ such that the head component converges strongly,
\[
a^{(n)}_{\mathrm{head}} \to a \quad \text{and} \quad b^{(n)}_{\mathrm{head}} \to b.
\]
\item The coordinates of the tail component vanishes,
\[
\max_{i > m_n} \max\bigl( |a_i^{(n)}|, |b_i^{(n)}| \bigr) \longrightarrow 0;
\]
\medskip

\item The norms of tail component converge
\[
\|a^{(n)}_{\mathrm{tail}}\|_2^2 \longrightarrow 1-\|a\|_2^2,
\qquad
\|b^{(n)}_{\mathrm{tail}}\|_2^2 \longrightarrow 1-\|b\|_2^2.
\]
In particular,
\[
\langle a^{(n)}_{\mathrm{tail}}, b^{(n)}_{\mathrm{tail}} \rangle
\longrightarrow
- \langle a,b \rangle.
\]
\end{enumerate}
\end{lemma}

\begin{proof}
After passing to a subsequence, we may first assume that
\[
\bigl(a_1^{(n)},b_1^{(n)}\bigr) \longrightarrow (a_1,b_1)
\]
for some scalars $a_1,b_1 \in \R$. Next, for each $n$ and $i \ge 2$, define
\[
s_i^{(n)} := \max\bigl( |a_i^{(n)}|, |b_i^{(n)}| \bigr),
\]
and permute the coordinates $\{2,3,\dots\}$ so that $(s_i^{(n)})_{i \ge 2}$ is
nonincreasing. By a diagonal extraction we may therefore assume that, for every fixed
$i \ge 2$,
\[
a_i^{(n)} \longrightarrow a_i,
\qquad
b_i^{(n)} \longrightarrow b_i,
\]
and define $a$ and $b$ as the limiting sequence
\[
a := (a_1,a_2,a_3,\dots),
\qquad
b := (b_1,b_2,b_3,\dots).
\]
We first verify that the limiting sequence is well-defined. In fact, by Fatou's lemma
\[
\|a\|_2^2=\sum_{i \ge 1} |a_i|^2
\le
\liminf_{n \to \infty} \sum_{i \ge 1} |a_i^{(n)}|^2
= 1,
\]
which implies that $a\in \ell_2(\mathbb{N})$. Similarly, we conclude that $b\in \ell^2(\mathbb{N})$.

Next, we construct the head component by first constructing the sequence $m_k$. To this end, notice that
\[
\sum_{i \ge 2} \bigl(s_i^{(n)}\bigr)^2
\le
\sum_{i \ge 2} \bigl( |a_i^{(n)}|^2 + |b_i^{(n)}|^2 \bigr)
\le 2,
\]
which together with the fact that $s_i^{(n)}$ is nonincreasing implies that 
\begin{equation}
\label{eq:uniformbound_si}
s_i^{(n)} \le \sqrt{\frac{2}{i-1}}
\qquad
\text{for every } n \ge 1,\ i \ge 2.
\end{equation}


Therefore, we may choose an increasing sequence of positive integers $(m_k)_{k \ge 1}$ such that
\[
\sum_{i > m_k} |a_i|^2 + \sum_{i > m_k} |b_i|^2 \le 2^{-k}
\qquad
\text{for every } k \ge 1.
\]
Consequently, for each fixed $k$ we have
\[
\sum_{i=1}^{m_k} |a_i^{(n)}-a_i|^2
+ \sum_{i=1}^{m_k} |b_i^{(n)}-b_i|^2
\longrightarrow 0
\qquad
\text{as } n \to \infty.
\]
Passing to a further subsequence in $n$, we may therefore assume that
\[
\sum_{i=1}^{m_n} |a_i^{(n)}-a_i|^2
+ \sum_{i=1}^{m_n} |b_i^{(n)}-b_i|^2
\le 2^{-n}
\qquad
\text{for every } n \ge 1.
\]
To see why $a^{(n)}_{\mathrm{head}}$ converges to $a$, observe that
\[
\|a^{(n)}_{\mathrm{head}}-a\|_2^2
\le
|a_1^{(n)}-a_1|^2
+ \sum_{i=2}^{m_n} |a_i^{(n)}-a_i|^2 + \sum_{i > m_n} |a_i|^2,
\]
which vanishes as $n$ goes to infinity. The same argument gives
$b^{(n)}_{\mathrm{head}} \to b$ in $\ell^2$. To verify item (II), the uniform estimate \eqref{eq:uniformbound_si} implies that
\[
\max_{i > m_n} \max\bigl( |a_i^{(n)}|, |b_i^{(n)}| \bigr)
\le
\sqrt{\frac{2}{m_n}},
\]
which vanishes as $m_n$ diverges. Finally, item (I) implies that
\[
\|a^{(n)}_{\mathrm{tail}}\|_2^2
=
1-\|a^{(n)}_{\mathrm{head}}\|_2^2
\longrightarrow
1-\|a\|_2^2,
\]
and similarly that $\|b^{(n)}_{\mathrm{tail}}\|_2^2$ converges to $1-\|b\|_2^2$. In particular, recalling that $a^{(n)}$ and $b^{(n)}$ are orthogonal
\[
0 = \langle a^{(n)}, b^{(n)} \rangle
=
\langle a^{(n)}_{\mathrm{head}}, b^{(n)}_{\mathrm{head}} \rangle
+
\langle a^{(n)}_{\mathrm{tail}}, b^{(n)}_{\mathrm{tail}} \rangle,
\]
and by item (I) again we have
\[
\langle a^{(n)}_{\mathrm{tail}}, b^{(n)}_{\mathrm{tail}} \rangle
\longrightarrow
- \langle a,b \rangle,
\]
which concludes the proof of item (III).
\end{proof}

\begin{remark}
The permutation in Lemma \ref{lem:profile} is applied only to the good
coordinates
\[
(\xi_{n,i},a_i^{(n)},b_i^{(n)})_{i \ge 2}.
\]
The exceptional coordinate $\eta$ is kept fixed at index $1$.
\end{remark}

\section{Auxiliary lemmas}

We now collect the three auxiliary ingredients used later to analyze the head
and the tail: bounded probes for recovering coefficients from the quadratic
observable, a uniform lower $L^1$-$L^2$ estimate for independent sums, and a
rigidity statement for diffuse tails.

\begin{proposition}[Bounded coefficient probes]
\label{prop:bounded-probes}
Let $\xi$ be a real
random variable satisfying for some $\delta \in (0,1)$
\[
\E[\xi] = 0,
\qquad
\|\xi\|_{L^2}=1
\qquad
\text{and}  \quad 
\delta \le \|\xi\|_{L^1} \le 1-\delta.
\]
Then there exist bounded random variables $S,T$ such that
\[
\E[S] = \E[T] = 0,
\qquad
\E[\xi S] = 1,
\qquad
\E[(\xi^2-1)T] = 1,
\]
and
\[
\|S\|_{L^\infty}
+ \|T\|_{L^\infty}
\le \frac{4}{\delta}.
\]
\end{proposition}

\begin{proof}
We first provide the construction of the random variable $S$. To this end, we define
\[
S_0 := \operatorname{sign}(\xi),
\]
and observe that
\[
\E[\xi S_0] = \E[|\xi|] = \|\xi\|_{L^1} \ge \delta.
\]
Thus the random variable
\[
S := \frac{S_0 - \E[S_0]}{\E[\xi S_0]}
\]
is well-defined and centered. Since $\E [\xi]=0$, we also have
\[
\E[\xi S] = \frac{\E [\xi S_0]-\E[\xi]\E[S_0]}{\E[\xi S_0]} = \frac{\E [\xi S_0]}{\E [\xi S_0]}=1.
\]
Moreover
\[
\|S\|_{L^\infty}
\le
\frac{2}{\E[|\xi|]}
\le
\frac{2}{\delta},
\]
which completes all the claims about $S$ for the constant $C_{\delta}=2/\delta\ge 2$. We now mimic the construction of $S$ to construct $T$ just replacing $\xi$ by $\xi^2-1$. More accurately, let 
\[
T_0 := \operatorname{sign}(\xi^2-1),
\]
we verify that $\E[(\xi^2-1)T_0]$ is bounded away from zero. Indeed,

\[
|\xi^2-1|
=
\bigl||\xi|-1\bigr| (|\xi|+1)
\ge
\bigl||\xi|-1\bigr|,
\]
so Jensen's inequality implies that
\[
\E[(\xi^2-1)T_0]=\E[|\xi^2-1|]
\ge
\E\big[\bigl||\xi|-1\bigr|\big]
\ge
\bigl| \E[|\xi|] - 1 \bigr|
=
1-\E[|\xi|]
\ge
\delta.
\]
The random variable
\[
T := \frac{T_0 - \E[T_0]}{\E[(\xi^2-1)T_0]}
\]
is centered and well-defined. Since $\E[\xi^2-1]=0$, we also have that 
\[
\E[(\xi^2-1)T] = \frac{\E[(\xi^2-1)T_0]-\E[(\xi^2-1)]\E T_0}{\E[(\xi^2-1)T_0]} = \frac{\E[(\xi^2-1)T_0]}{\E[(\xi^2-1)T_0]}=1.
\]
Finally, 
\[
\|T\|_{L^\infty}
\le
\frac{2}{\E[|\xi^2-1|]}
\le
\frac{2}{\delta}.
\]
\end{proof}

\begin{remark}
For the compactness argument, Proposition \ref{prop:bounded-probes} is exactly
what is needed in order to recover the coefficients $a_i^2$ and $a_i a_j$ from
the quadratic observable $|X|^2$.
\end{remark}

\begin{proposition}[Probes for the exceptional coordinate]
\label{prop:exceptional-probes}
Let $\eta$ be a real-valued random variable satisfying
\[
\E[\eta]=0, \quad \|\eta\|_{L^2}=1.
\]
Then there exists a bounded random variable $S_\eta$ such that
\[
\E[S_\eta]=0,
\qquad
\E[\eta S_\eta]=1.
\]
If $\eta^2 \not \equiv  1$ almost surely, then there exists a bounded random variable
$T_\eta$ such that
\[
\E[T_\eta]=0,
\qquad
\E[(\eta^2-1)T_\eta]=1.
\]
\end{proposition}

\begin{proof}
Let
\[
S_0 := \operatorname{sign}(\eta).
\]
Since $\|\eta\|_{L^2}=1$, one has $\E[|\eta|]>0$. Thus
\[
S_\eta := \frac{S_0-\E[S_0]}{\E[\eta S_0]}
\]
is bounded and satisfies
\[
\E[S_\eta]=0,
\qquad
\E[\eta S_\eta]=1.
\]
Assume now that $\eta^2 \not \equiv 1$ almost surely. Let
\[
T_0 := \operatorname{sign}(\eta^2-1).
\]
Since $\E[|\,\eta^2-1\,|]>0$, the bounded function
\[
T_\eta := \frac{T_0-\E[T_0]}{\E[(\eta^2-1)T_0]}
\]
satisfies
\[
\E[T_\eta]=0,
\qquad
\E[(\eta^2-1)T_\eta]=1.
\]
\end{proof}
The next lemma is a standard combination of a symmetrization argument and Khintchine inequality. We opt for a self-contained statement.

\begin{lemma}[Lower $L^1$-$L^2$ comparability is stable under independent sums]
\label{lem:l1l2-linear-combinations}
Let $(\zeta_i)_{i=1}^N$ be independent centered real-valued random variables and
assume that
\[
\|\zeta_i\|_{L^1} \ge \delta \|\zeta_i\|_{L^2}
\qquad
\text{for every } i=1,\dots,N.
\]
Then for every scalar family $(c_i)_{i=1}^N$ one has
\[
\left\|\sum_{i=1}^N c_i \zeta_i \right\|_{L^1}
\ge
\frac{\delta}{2\sqrt 2}
\left(\sum_{i=1}^N c_i^2 \|\zeta_i\|_{L^2}^2 \right)^{1/2}.
\]
In particular, if $\|\zeta_i\|_{L^2}=1$ for all $i$, then
\[
\left\|\sum_{i=1}^N c_i \zeta_i \right\|_{L^1}
\ge
\frac{\delta}{2\sqrt 2}\|(c_i)_{i=1}^N\|_{\ell^2}.
\]
\end{lemma}

\begin{proof}
By homogeneity, it is enough to consider the  following case
\[
\sum_{i=1}^N c_i^2 \|\zeta_i\|_{L^2}^2 = 1.
\]
By replacing the random variables $\zeta_i$ by $\zeta_i/\|\zeta_i\|_{L^2}$, we may without loss of generality assume
that $\|\zeta_i\|_{L^2}=1$ and consequently, $\sum_{i=1}^N c_i^2=1$.

Set 
\[
S := \sum_{i=1}^N c_i \zeta_i,
\]
and consider an independent copy of $S'$ given by
\[
S' := \sum_{i=1}^N c_i \zeta_i'.
\]
We consider a symmetrized version of $S$. Indeed, notice that
\[
2\|S\|_{L^1}
\ge
\|S-S'\|_{L^1}
=
\left\|\sum_{i=1}^N c_i(\zeta_i-\zeta_i')\right\|_{L^1}.
\]
Now the differences $\zeta_i-\zeta_i'$ are independent and symmetric random variables. Equivalently, for an independent Rademacher random variable $r_i$, the random variables $\zeta_i-\zeta_i'$ and $r_i(\zeta_i-\zeta_i')$ are equal in distribution. Hence, by conditioning on the differences $\zeta_i-\zeta_i'$ we apply Khintchine's inequality to the $r_i$'s to obtain that
\[
\left\|\sum_{i=1}^N c_i(\zeta_i-\zeta_i')\right\|_{L^1} = \left\|\sum_{i=1}^N c_ir_i(\zeta_i-\zeta_i')\right\|_{L^1}
\ge
\frac{1}{\sqrt 2}
\E\left[\left(\sum_{i=1}^N c_i^2(\zeta_i-\zeta_i')^2\right)^{1/2}\right].
\]
Recall that we assume that $\sum_i c_i^2=1$. Thus, Cauchy--Schwarz yields the pointwise bound
\[
\sum_{i=1}^N c_i^2 |\zeta_i-\zeta_i'| = \sum_{i=1}^N c_i c_i |\zeta_i-\zeta_i'|\le 
\underbrace{\left(\sum_{i=1}^N c_i^2\right)}_{=1}\left(\sum_{i=1}^N c_i^2(\zeta_i-\zeta_i')^2\right)^{1/2},
\]
which implies that
\[
\|S\|_{L^1}
\ge
\frac{1}{2\sqrt 2}
\sum_{i=1}^N c_i^2 \E[|\zeta_i-\zeta_i'|].
\]
To conclude the proof of the theorem, it is enough to prove that $\E[|\zeta_i-\zeta_i'|]\ge \delta$ because if this is the case then
\[
\|S\|_{L^1}
\ge
\frac{1}{2\sqrt 2}
\sum_{i=1}^N c_i^2 \delta = \frac{\delta}{2\sqrt{2}},
\]
which is exactly the normalized version of the conclusion of the theorem. To see why this is true, recall that
by the independence between $\zeta_i$ and $\zeta_i'$, thus Jensen's inequality implies that
\[
\E[\zeta_i-\zeta_i'|]
\ge
\E\big|\E[\zeta_i-\zeta_i' \mid \zeta_i]\big|
=
\E[|\xi_i|]
\ge
\delta,
\]
where $\E[\zeta_i-\zeta_i'|\zeta_i]=\zeta_i$ follows from the fact that $\zeta_i'$ is centered and its law does not change once conditioned on $\zeta_i$. 
\end{proof}

\textcolor{red}{Pedro: I commented the second assumption on $c^{(n)}$, I think it is not needed here}
\begin{proposition}[Rigidity of one-dimensional tail limits]
\label{prop:tail-rigidity}
Let $(\xi_{n,i})_{n \ge 1,\ i \ge 2}$ be a triangular array as in
Proposition \ref{prop:counterexample-array}, and let $c^{(n)} \in \ell^2(\N)$
be a sequence of vectors satisfying
\[
\sup_{n\ge1} \|c^{(n)}\|_2 < \infty.
\]
%and
%\[
%\max_{i \ge 2} |c_i^{(n)}| \longrightarrow 0,
%\qquad
%c_1^{(n)} = 0 \text{ for every } n.
%\]
If the sequence
\[
Z_n := \sum_{i \ge 2} c_i^{(n)} \xi_{n,i}
\]
converges to zero in distribution, then the sequence $\|c^{(n)}\|_2$ vanishes. Equivalently, if along a subsequence the random vector
\[
\left(
\sum_{i \ge 2} a_i^{(n)} \xi_{n,i},
\sum_{i \ge 2} b_i^{(n)} \xi_{n,i}
\right)
\]
converges in distribution to a random vector supported on the line $\alpha u + \beta v = 0$ for some
$(\alpha,\beta) \ne (0,0)$, then along that same subsequence,
\[
\|\alpha a^{(n)} + \beta b^{(n)}\|_2 \longrightarrow 0.
\]
\end{proposition}

\begin{proof}
Suppose the sequence $Z_n$ converges in distribution and the limit is the constant zero. Since
convergence in distribution to a constant implies convergence in probability, we have that $Z_n$ converges in probability to zero. On the other hand,
\[
\|Z_n\|_{L^2} = \|c^{(n)}\|_2
\]
is uniformly bounded by the first assumption on $c^{(n)}$. Thus the family $(|Z_n|)_n$ is uniformly
integrable, and by Vitali's convergence theorem
\[
\|Z_n\|_{L^1} \longrightarrow 0.
\]
Applying Lemma \ref{lem:l1l2-linear-combinations} row by row to the family
$(\xi_{n,i})_{i \ge 2}$ gives
\[
\|Z_n\|_{L^1}
\ge
\frac{\delta}{2\sqrt2}\|c^{(n)}\|_2.
\]
Hence the sequence $\|c^{(n)}\|_2$ vanishes. For the equivalent formulation, we apply the first part to
\[
c_i^{(n)} := \alpha a_i^{(n)} + \beta b_i^{(n)},
\]
which is a valid choice as for every $n\ge 1$
\[
\|c^{(n)}\|_2
\le
|\alpha|\|a^{(n)}\|_2 + |\beta|\|b^{(n)}\|_2
=
|\alpha|+|\beta|.
\]
\end{proof}

\section{Identification of the limit}
In this section we complete the preliminaries results to prove our main theorem. With the head-tail decomposition and the auxiliary estimates in hand, the first step is to
identify the limiting head profile from the fact that the quadratic information
$|X_n|^2-|Y_n|^2$ converges strongly to zero. After that, we exploit rigidity of the weak limit of full head-tail
decomposition.

\begin{lemma}[Identification of the head profile]\label{lem:head-sign}
Consider the head and tail decomposition from Proposition \ref{prop:counterexample-array} and Lemma \ref{lem:profile}. Then one of the following alternatives must be true:
either there exists $\tau \in \{\pm 1\}$ such that
\[
a = \tau b
\qquad
\text{and}
\qquad
\|a^{(n)}_{\mathrm{head}} - \tau b^{(n)}_{\mathrm{head}}\|_2 \longrightarrow 0,
\]
or  $\eta$ is a Rademacher variable and
\[
a_i=b_i=0
\qquad
\text{for every } i \ge 2.
\]
\end{lemma}

\begin{proof}
Recall that Proposition \ref{prop:counterexample-array} asserts that the random variable $|X_n|-|Y_n|$ converges to zero in $L^2$. Hence by Cauchy-Schwarz inequality,
\begin{equation}
\label{eq:L1_quadratic_vanishes}
\bigl\| |X_n|^2 - |Y_n|^2 \bigr\|_{L^1}
\le
\|\, |X_n|-|Y_n| \,\|_{L^2}
\cdot
\|\, |X_n|+|Y_n| \,\|_{L^2}
\le
2 \|\, |X_n|-|Y_n| \,\|_{L^2}
\longrightarrow 0.
\end{equation}
Fix an index $i \ge 2$ and apply Proposition \ref{prop:bounded-probes} to $\xi_{n,i}$ to obtain a probe $T_{n,i}$ that satisfies $\|T_{n,i}\|_{L^\infty} \le
C_\delta$, for some positive constant $C_{\delta}$. Thus \eqref{eq:L1_quadratic_vanishes} implies that
\begin{equation}
\label{eq:quadraticxprobe_vanishes}
\E\bigl[(|X_n|^2-|Y_n|^2) T_{n,i}\bigr] \le \bigl\| |X_n|^2 - |Y_n|^2 \bigr\|_{L^1}C_{\delta} \longrightarrow 0 .
\end{equation}
On the other hand, recalling the definition of $X_n$ in Proposition \ref{prop:counterexample-array} we have
\begin{equation*}
\E\bigl[|X_n|^2 T_{n,i}\bigr] = (a_1^{(n)})^2\E[T_{n,i}] + 2a_1^{(n)}\E[\eta] \E\big[\sum_{j\ge 2}a_j^{(n)}\xi_{n,j}T_{n,i}\big] + \E\big[\sum_{j,j'\ge 2}a_j^{(n)}a_{j'}^{(n)}\xi_{n,j}\xi_{n,j'}T_{n,i}\big].
\end{equation*}
Since $T_{n,i}$ and $\eta$ are centered and independent the first two terms on the right-hand side are zero. Moreover, note that $T_{n,i}$ is a random variable that depends only on $\xi_{n,i}$ alone, so it is independent of
$(\eta,\xi_{n,j})_{j \ne i}$ which implies that 
\begin{equation*}
\E\bigl[|X_n|^2 T_{n,i}\bigr] =  \E\big[\sum_{j,j'\ge 2}a_j^{(n)}a_{j'}^{(n)}\xi_{n,j}\xi_{n,j'}T_{n,i}\big] = (a_i^{(n)})^2\E[\xi_{n,i}^2T_{n,i}] =(a_i^{(n)})^2,
\end{equation*}
where the last step follows Proposition \ref{prop:bounded-probes}. Similarly, we obtain that
\[
\E\bigl[(|X_n|^2-|Y_n|^2) T_{n,i}\bigr]
=
(a_i^{(n)})^2 - (b_i^{(n)})^2,
\]
hence together with \eqref{eq:quadraticxprobe_vanishes},
\[
(a_i^{(n)})^2 - (b_i^{(n)})^2 \longrightarrow 0.
\]
Passing to the limit gives
\begin{equation}
\label{eq:squared_good}
a_i^2 = b_i^2,
\qquad
\text{for every } i \ge 2.
\end{equation}


Now we fix a pair of indexes $i,j \ge 2$ with $i \ne j$. Let $S_{n,i},S_{n,j}$ be the corresponding probes given by
Proposition \ref{prop:bounded-probes} for $\xi_{n,i}$ and $\xi_{n,j}$, respectively.
By the boundness of the probes and $L^1$ convergence from \eqref{eq:L1_quadratic_vanishes}
\[
\E\bigl[(|X_n|^2-|Y_n|^2) S_{n,i} S_{n,j}\bigr] \le \|S_{n,i}\|_{L^{\infty}}\|S_{n,j}\|_{L^{\infty}} \||X_n|^2-|Y_n|^2\|_{L^1}\longrightarrow 0.
\]
Expanding $X_n^2$ as before and invoking properties of $S_{n,i}$, $S_{n,j}$ from Proposition \ref{prop:bounded-probes},
\[
\E[X_n^2 S_{n,i}S_{n,j}] = 2a_i^{(n)}a_j^{(n)}\E[\xi_{n,i}\xi_{n,j}S_{n,i}S_{n,j}] = 2a_i^{(n)}a_j^{(n)}\E[\xi_{n,i}S_{n,i}]\E[\xi_{n,j}S_{n,j}] = 2a_i^{(n)}a_j^{(n)}.
\]
Similarly,
\[
\E\bigl[(|X_n|^2-|Y_n|^2) S_{n,i} S_{n,j}\bigr]
= 2a_i^{(n)}a_j^{(n)} - 2b_i^{(n)}b_j^{(n)},
\]
therefore in the limit
\begin{equation}
\label{eq:crossterms_good}
a_i a_j = b_i b_j, \qquad
\text{for every } i,j \ge 2,\ i \ne j.
\end{equation}


We now repeat the argument for $\eta$ by
invoking Proposition \ref{prop:exceptional-probes} to obtain the probe $S_{\eta}$. For every
$j \ge 2$, let $S_{n,j}$ be the probe given by Proposition
\ref{prop:bounded-probes} for $\xi_{n,j}$. By repeating the same argument, we obtain
\[
\E\bigl[(|X_n|^2-|Y_n|^2) S_\eta S_{n,j}\bigr] \longrightarrow 0,
\]
and
\[
\E\bigl[(|X_n|^2-|Y_n|^2) S_\eta S_{n,j}\bigr]
=
2\bigl(a_1^{(n)} a_j^{(n)} - b_1^{(n)} b_j^{(n)}\bigr).
\]
We conclude that
\[
a_1 a_j = b_1 b_j
\qquad
\text{for every } j \ge 2.
\]
\textbf{First Alternative Case 1:}
Assume first that there exists $i_0 \ge 2$ with $a_{i_0} \ne 0$. Since
$a_{i_0}^2=b_{i_0}^2$ by \eqref{eq:squared_good}, we have $b_{i_0} = \tau a_{i_0}$ for some
$\tau \in \{\pm 1\}$. For every $j \ge 2$, \eqref{eq:crossterms_good} implies that
\[
a_{i_0} a_j = b_{i_0} b_j = \tau a_{i_0} b_j,
\]
so $a_j = \tau b_j$. In particular, for  $j=i_0$ we have
\[
a_1 a_{i_0} = b_1 b_{i_0} = \tau a_{i_0} b_1,
\]
which implies that $a_1 = \tau b_1$. We conclude that $a=\tau b$ for some $\tau \in \{\pm 1\}$.

\textbf{First Alternative Case 2:}
Assume now that $a_i=0$ for every $i\ge 2$ and $\eta^2 \not \equiv 1$ almost surely. Then $b_i=0$ for every $i \ge 2$ by \eqref{eq:squared_good}. Since we assume that $\eta^2 \not\equiv 1$ almost
surely, let $T_\eta$ be the probe given by Proposition
\ref{prop:exceptional-probes}. Again using boundedness and $L^1$ convergence,
\[
\E\bigl[(|X_n|^2-|Y_n|^2) T_\eta\bigr] \longrightarrow 0,
\]
and by independence
\[
\E\bigl[(|X_n|^2-|Y_n|^2) T_\eta\bigr]
=
(a_1^{(n)})^2-(b_1^{(n)})^2.
\]
Hence $a_1^2=b_1^2$, so $a=\tau b$ for some $\tau \in \{\pm 1\}$. Therefore if $a_i>0$ for some $i\ge 2$ or $a_i=0$ for every $i\ge 2$ and $\eta^2 \not \equiv1$ a.s we have 
\[
\|a^{(n)}_{\mathrm{head}} - \tau b^{(n)}_{\mathrm{head}}\|_2
\le
\|a^{(n)}_{\mathrm{head}}-a\|_2
+
\|b^{(n)}_{\mathrm{head}}-b\|_2
\longrightarrow 0.
\]
\textbf{Second Alternative:} The only case left is when $a_i=0$ for every $i\ge 2$ and $\eta^2\equiv1$ almost surely. Since $\E[\eta]=0$ and $\|\eta\|_{L^2}=1$, we must have that $\eta$ is a Rademacher variable.
This is exactly the second alternative.


\end{proof}

\begin{lemma}[Limit law for the head-tail decomposition]
\label{lem:limit-law}
Under the notation from Lemma \ref{lem:profile}, let
\[
X_n^{\mathrm{head}} := a_1^{(n)}\eta + \sum_{2 \le i \le m_n} a_i^{(n)} \xi_{n,i},
\qquad
X_n^{\mathrm{tail}} := \sum_{i > m_n} a_i^{(n)} \xi_{n,i},
\]
and define the random variables $Y_n^{\mathrm{head}}$ and $Y_n^{\mathrm{tail}}$ analogously. Then there exists a subsequence for which the random vector 
\[
\bigl(
X_n^{\mathrm{head}},
Y_n^{\mathrm{head}},
X_n^{\mathrm{tail}},
Y_n^{\mathrm{tail}}
\bigr)
\]
converges in distribution to a random vector
\[
\bigl(
H_X,H_Y,R_X,R_Y
\bigr)
\in \R^4,
\]
where:
\begin{enumerate}
\renewcommand{\labelenumi}{\Roman{enumi}.}
\item The random vector $(H_X,H_Y)$ is independent of $(R_X,R_Y)$;
\item The law of $(R_X,R_Y)$ is infinitely divisible;
\item One of the following alternatives must be true: there exist s$\tau \in \{\pm 1\}$ for which 
\[
H_X = \tau H_Y, \quad \text{almost surely.}
\]
Or $\eta$ is a Rademacher variable and
\[
H_X = a_1 \eta,
\qquad
H_Y = b_1 \eta
\qquad
\text{almost surely;}
\]
\item It holds that
\[
|H_X + R_X| = |H_Y + R_Y|
\qquad
\text{almost surely.}
\]
\end{enumerate}
\end{lemma}

\begin{proof}
Each of the entries of $\bigl(
X_n^{\mathrm{head}},
Y_n^{\mathrm{head}},
X_n^{\mathrm{tail}},
Y_n^{\mathrm{tail}}
\bigr)$ is bounded in $L^2$ by one, so the family of
quadruples is tight. We apply Prohorov's theorem to obtain that the existence of a limit $(H_X,H_Y,R_X,R_Y)$ in distribution.

For every $n$, the pair
$(X_n^{\mathrm{head}},Y_n^{\mathrm{head}})$ depends only on $\eta$ and on
$(\xi_{n,i})_{2 \le i \le m_n}$, while
$(X_n^{\mathrm{tail}},Y_n^{\mathrm{tail}})$ depends only on
$(\xi_{n,i})_{i > m_n}$. Thus the pairs are independent and since independence is preserved under convergence in distribution, we obtain that
$(H_X,H_Y)$ is independent of $(R_X,R_Y)$.

To verify item (II), consider the row-wise array of
summands
\[
Z_{n,i}
:=
\bigl(a_i^{(n)} \xi_{n,i},\, b_i^{(n)} \xi_{n,i}\bigr)\mathbf{1}_{\{i>m_n\}}.
\]
The array is infinitesimal: for every $\varepsilon>0$, Chebyshev's inequality combined with item (II) from Lemma \ref{lem:profile} implies that
\[
\sup_{i>m_n}
\mathbb P\bigl(|Z_{n,i}|>\varepsilon\bigr)
\le
\frac{1}{\varepsilon^2}
\sup_{i>m_n}\bigl((a_i^{(n)})^2+(b_i^{(n)})^2\bigr)
\longrightarrow 0.
\]
For each fixed $n$, the summands $(Z_{n,i})_{i > m_n}$ are independent as the random variables $\xi_{n,i}$ are independent. By the classical theorem on limits of
row-wise independent infinitesimal triangular arrays
\cite[Chapter~4]{gnedenko-kolmogorov}, every limit of the
sums $\sum_i Z_{n,i}$ is infinitely divisible.

Next, we verify item (III). To this end, suppose the first alternative of Lemma \ref{lem:head-sign}. Then
\[
\|X_n^{\mathrm{head}} - \tau Y_n^{\mathrm{head}}\|_{L^2}
=
\|a^{(n)}_{\mathrm{head}} - \tau b^{(n)}_{\mathrm{head}}\|_2
\longrightarrow 0,
\]
hence $H_X = \tau H_Y$ almost surely. On the other hand, if we are in the second alternative of Lemma \ref{lem:head-sign}, then
\[
\sum_{2 \le i \le m_n} (a_i^{(n)})^2
\longrightarrow 0,
\qquad
\sum_{2 \le i \le m_n} (b_i^{(n)})^2
\longrightarrow 0,
\]
and therefore
\[
\|X_n^{\mathrm{head}} - a_1 \eta\|_{L^2}
\le
|a_1^{(n)}-a_1|
+
\left(\sum_{2 \le i \le m_n} (a_i^{(n)})^2\right)^{1/2}
\longrightarrow 0,
\]
and the same argument implies that $Y_n^{\mathrm{head}}$ converges to $b_1\eta$ in $L^2$. Thus
\[
H_X = a_1 \eta,
\qquad
H_Y = b_1 \eta
\qquad
\text{almost surely,}
\]
which concludes the second alternative in item (III). Finally, to verify item (IV), we observe that
\[
\bigl|X_n^{\mathrm{head}} + X_n^{\mathrm{tail}}\bigr|
-
\bigl|Y_n^{\mathrm{head}} + Y_n^{\mathrm{tail}}\bigr|
=
|X_n|-|Y_n|
\longrightarrow 0
\]
in $L^2$, hence in probability. On the other hand, by the continuous mapping
theorem this difference also converges in distribution to
$|H_X+R_X|-|H_Y+R_Y|$. By uniqueness of the limit, the random variable $|H_X+R_X|-|H_Y+R_Y|$ must be equal to zero in distribution which implies that it is zero almost surely.
\end{proof}
\textcolor{blue}{Pedro: I think we only use for lines passing though the origin, in this case the proof is simpler. Or reduce to this case by shifting...}
\begin{lemma}[A geometric two-line support lemma]
\label{lem:two-line}
Let $\mu$ be an infinitely divisible probability measure on $\R^2$. If
$\supp \mu$ is contained in the union of two affine lines, then
$\supp \mu$ is contained in a single affine line.
\end{lemma}


\begin{proof}
Because $\mu$ is infinitely divisible, it is in particular $2$-divisible:
there exists a probability measure $\nu$ such that $\mu=\nu * \nu$. Let $A$ be the support of the measure $\nu$. We claim that the set $A$ is contained in a single affine line. If such claim is true it follows from
\[
\supp \mu = \overline{A+A},
\]
that $\supp(\mu)$ is contained in a single affine line.

To prove the claim, assume by contradiction that $A$ is not contained in a single affine line. Then there exist three points $x,y,z \in A$ that are not collinear. By assumption,
\[
\supp \mu \subset L_1 \cup L_2,
\]
where $L_1$ and $L_2$ are two affine lines which we assume that they are distinct, otherwise the claim is trivial.
Hence, the six points
\[
2x,\ 2y,\ 2z,\ x+y,\ x+z,\ y+z
\]
all belong to $A+A\subset L_1 \cup L_2$. Among $2x,2y,2z$, two must lie on the same line; without loss of generality, assume that $2x,2y \in L_1$. In other words, $L_1$ is the affine line generated by $x$ and $y$.

Since $x,y,z$ are not collinear, the points $2z, x+z, y+z$ cannot belong to $L_1$. Hence the only option is that
\[
2z,\ x+z,\ y+z \in L_2.
\]
However this implies that $x-y=(x+z)-(y+z)$ is in the direction of $L_2$ and  also in the direction of $L_1$. Consequently, $L_1$ and $L_2$ must be parallel and also as $2x=2(x+z)-2z$ we have that the lines intersect. We conclude that $L_1$ and $L_2$ coincide which contradicts our assumption that they are distinct. This proves the claim that $A$ is contained in a single affine line, and so is $\supp(\mu)$.


%three non-collinear points $x,y,z \in A$. Let $L_1 \cup L_2$ be the union of
%two affine lines containing $\supp \mu$. Since $A+A \subset \supp \mu$, the six
%points
%\[
%2x,\ 2y,\ 2z,\ x+y,\ x+z,\ y+z
%\]
%belong to $L_1 \cup L_2$. Among the three points $2x,2y,2z$, two must belong
%to the same line; after relabeling, assume $2x,2y \in L_1$.

%Then $L_1$ is the line through $2x$ and $2y$. If $x+z \in L_1$, then $x,y,z$
%would be collinear, which is impossible. The same argument shows that
%$y+z \notin L_1$ and $2z \notin L_1$. Consequently,
%\[
%2z,\ x+z,\ y+z \in L_2.
%\]
%But $x+z$ and $y+z$ lie on the line $z+\R(y-x)$, whereas $2z$ belongs to that
%line only if $z,x,y$ are collinear. This contradiction proves that $A$ is
%contained in an affine line, and %therefore so is
%$\supp \mu = \overline{A+A}$.
\end{proof}

\section{Proof of Theorem \ref{thm:main}}

We are now ready to complete the contradiction argument by combining the structural
information obtained for the strong limit for the head component with the rigidity of the limit in distribution of the tail component.

\begin{proof}
Assume for the sake of contradiction that Theorem \ref{thm:main} is false, and consider the counterexample array be
given by Proposition \ref{prop:counterexample-array}. We follow the notation introduced in the 
Lemmata \ref{lem:profile}, \ref{lem:head-sign} and \ref{lem:limit-law}.

By Lemma \ref{lem:limit-law} item (I), we know that $(H_X,H_Y)$ is independent of $(R_X,R_Y)$, thus we may realize the joint law of the pairs
on a product space. In particular, there exists a set $\Omega_0$ of probability
one with respect to the head variables $(H_X,H_Y)$ such that, for every $\omega \in \Omega_0$,
\begin{equation}
\label{eq:a.s_equality_tail}
|H_X(\omega)+R_X| = |H_Y(\omega)+R_Y|
\qquad
\text{holds almost surely in the tail variables.}
\end{equation}

\textbf{Second Alternative of Lemma
\ref{lem:head-sign}.}
In this case $\eta$ is a Rademacher variable and, by Lemma
\ref{lem:limit-law} item (III),
\[
H_X = a_1 \eta,
\qquad
H_Y = b_1 \eta
\qquad
\text{almost surely.}
\]
Choose a realization $\omega_+,\omega_- \in \Omega_0$ satisfying
\[
\bigl(H_X(\omega_+),H_Y(\omega_+)\bigr)=(a_1,b_1),
\qquad
\bigl(H_X(\omega_-),H_Y(\omega_-)\bigr)=(-a_1,-b_1).
\]
Then it follows from \eqref{eq:a.s_equality_tail} that the support of the law of $(R_X,R_Y)$ is contained in the union of affine lines, that is, it is contained in both
\[
\Gamma_+
:=
\{(u,v)\in\R^2 : u-v=b_1-a_1\}
\cup
\{(u,v)\in\R^2 : u+v=-(a_1+b_1)\}
\]
and
\[
\Gamma_-
:=
\{(u,v)\in\R^2 : u-v=a_1-b_1\}
\cup
\{(u,v)\in\R^2 : u+v=a_1+b_1\}.
\]
By Lemma \ref{lem:limit-law}, we know that the law of $(R_X,R_Y)$ is infinitely divisible, therefore by Lemma \ref{lem:two-line}, the support of the law of $(R_X,R_Y)$ is
contained in a single affine line. Since $(R_X,R_Y)$ is centered (as a weak
limit of centered random vectors), the barycenter of the law lies at the
origin, so the line must pass through the origin.

\textbf{Case 1:} if $(R_X,R_Y)=(0,0)$ almost surely, then Proposition
\ref{prop:tail-rigidity} applied separately to
\[
\sum_{i > m_n} a_i^{(n)} \xi_{n,i}
\qquad
\text{and}
\qquad
\sum_{i > m_n} b_i^{(n)} \xi_{n,i}
\]
gives
\[
\|a^{(n)}_{\mathrm{tail}}\|_2 + \|b^{(n)}_{\mathrm{tail}}\|_2 \longrightarrow 0.
\]
Recalling that $a_i=b_i=0$ for every $i \ge 2$ by Lemma \ref{lem:head-sign}, we apply Lemma \ref{lem:profile} to obtain that 
\[
\|a^{(n)}_{\mathrm{head}}\|_2 \longrightarrow |a_1|,
\qquad
\|b^{(n)}_{\mathrm{head}}\|_2 \longrightarrow |b_1|.
\]
Recall that  $a^{(n)}$ and $b^{(n)}$ are unit norm, it forces the identity $|a_1|=|b_1|=1$. On the other hand, by construction $a^{(n)}$ and $b^{(n)}$ are orthogonal, so
\[
0
=
\langle a^{(n)},b^{(n)}\rangle
=
\langle a^{(n)}_{\mathrm{head}},b^{(n)}_{\mathrm{head}}\rangle
+
\langle a^{(n)}_{\mathrm{tail}},b^{(n)}_{\mathrm{tail}}\rangle,
\]
and the second term vanishes as the  norms vanish. Taking the limit, we get that gives
$a_1 b_1 = 0$ which contradicts $|a_1b_1|=1$.

\textbf{Case 2:} if $(R_X,R_Y)\not \equiv 0$ almost surely, then we may choose a nonzero point in the support of $(R_X,R_Y)$. Since the
support is contained in a line through the origin and in both $\Gamma_+$ and
$\Gamma_-$, that line must be the one with $u=v$ or $u=-v$.

If the support is contained in $u=v$, choose a nonzero point $(u,u)$ in the
support. Thus $(u,u) \in \Gamma_+ \cap \Gamma_-$ and one must have $a_1=b_1$.
Similarly, if the support is contained in $u=-v$, one must have $a_1=-b_1$. Thus there exists $\sigma \in \{\pm 1\}$ such that
\[
R_X=\sigma R_Y \text{ almost surely}
\qquad
\text{and}
\qquad
a_1=\sigma b_1.
\]
Consequently, Proposition \ref{prop:tail-rigidity} yields
\[
\|a^{(n)}_{\mathrm{tail}} - \sigma b^{(n)}_{\mathrm{tail}}\|_2 \longrightarrow 0,
\]
and $a=(a_1,0,0,\dots)$ and $b=(b_1,0,0,\dots)$ with $a_1=\sigma b_1$. By Lemma \ref{lem:profile} item (I), we have
\[
\|a^{(n)}_{\mathrm{head}} - \sigma b^{(n)}_{\mathrm{head}}\|_2 \longrightarrow 0.
\]
Hence
\[
\|a^{(n)}-\sigma b^{(n)}\|_2 \longrightarrow 0,
\]
again contradicting the orthonormality
\[
\|a^{(n)}-\sigma b^{(n)}\|_2^2 = 2.
\]
\textbf{First Alternative of Lemma \ref{lem:head-sign}.} In this case, we know that there exists
$\tau \in \{\pm 1\}$ such that
\[
H_X=\tau H_Y \text{ almost surely.}
\]
Then for every $\omega \in \Omega_0$ the support of the law of $(R_X,R_Y)$ is
contained in
\[
\Gamma_\omega :=
\begin{cases}
\{(u,v)\in\R^2 : u-v=0\} \cup \{(u,v)\in\R^2 : u+v=-2H_X(\omega)\},
& \text{if $\tau=1$}, \\[2mm]
\{(u,v)\in\R^2 : u+v=0\} \cup \{(u,v)\in\R^2 : u-v=-2H_X(\omega)\},
& \text{if $\tau=-1$}.
\end{cases}
\]
Repeating the same arguments, we know that the support of the law of $(R_X,R_Y)$ is
contained in a single affine line by Lemma \ref{lem:two-line} and that line
must pass through the origin.


\medskip
\noindent
\textbf{Case 1: $H_X \not\equiv 0$.}
Choose $\omega \in \Omega_0$ with $H_X(\omega)\ne 0$. Then among the two lines
forming $\Gamma_\omega$, the only one passing through the origin is either
\[
\{u-v=0\} \quad \text{if } \tau=1,
\qquad \text{or} \quad
\{u+v=0\} \quad \text{if } \tau=-1.
\]
Therefore the law of $(R_X,R_Y)$ is supported on the line $u=\tau v$. By
Proposition \ref{prop:tail-rigidity},
\[
\|a^{(n)}_{\mathrm{tail}} - \tau b^{(n)}_{\mathrm{tail}}\|_2 \longrightarrow 0,
\]
and by Lemma \ref{lem:head-sign} (first alternative)
\[
\|a^{(n)}_{\mathrm{head}} - \tau b^{(n)}_{\mathrm{head}}\|_2 \longrightarrow 0.
\]
Hence
\[
\|a^{(n)}-\tau b^{(n)}\|_2 \longrightarrow 0,
\]
which contradictis
\[
\|a^{(n)}-\tau b^{(n)}\|_2^2
= 2.
\]

\medskip
\noindent
\textbf{Case 2: $H_X \equiv 0$.}
Since $X_n^{\mathrm{head}}$ converges to $H_X$ in distribution, it must converge to zero in probability as well. In addition, $X_n^{\mathrm{head}}$ is uniformly integrate, thus Vitali's convergence theorem implies that
$\|X_n^{\mathrm{head}}\|_{L^1}$ vanishes. Let $S_\eta$ be the bounded random variable given by Proposition \ref{prop:exceptional-probes}, hence
\[
a_1^{(n)}
=
\E[X_n^{\mathrm{head}} S_\eta]\le \|S_\eta\|_{L^\infty}
\|X_n^{\mathrm{head}}\|_{L^1}
\longrightarrow 0,
\]
For every fixed $i \ge 2$, one has $i \le m_n$ for all sufficiently large $n$.
Let $S_{n,i}$ be the probe given by Proposition \ref{prop:bounded-probes} for
$\xi_{n,i}$. Then
\[
a_i^{(n)}
=
\E[X_n^{\mathrm{head}} S_{n,i}]
\le C_\delta
\|X_n^{\mathrm{head}}\|_{L^1} \longrightarrow 0,
\]
Passing to
the limit, we obtain that $a_i=0$ for every $i \ge 1$, hence $a=0$. Similarly, the same argument gives $b=0$, so Lemma \ref{lem:profile} yields
\[
\|a^{(n)}_{\mathrm{head}}\|_2 + \|b^{(n)}_{\mathrm{head}}\|_2 \longrightarrow 0.
\]
Since $H_X = \tau H_Y = 0$, the set $\Gamma_\omega$ reduces to
$\{u=v\} \cup \{u=-v\}$ for every $\omega$. By Lemma \ref{lem:two-line}, the
support of the law of $(R_X,R_Y)$ is contained in one of the two lines $u=v$
or $u=-v$. Thus there exists $\sigma \in \{\pm 1\}$ such that
$R_X=\sigma R_Y$ almost surely, and Proposition
\ref{prop:tail-rigidity} yields
\[
\|a^{(n)}_{\mathrm{tail}} - \sigma b^{(n)}_{\mathrm{tail}}\|_2 \longrightarrow 0.
\]
Therefore
\[
\|a^{(n)}-\sigma b^{(n)}\|_2
\le
\|a^{(n)}_{\mathrm{head}}\|_2
+
\|b^{(n)}_{\mathrm{head}}\|_2
+
\|a^{(n)}_{\mathrm{tail}}-\sigma b^{(n)}_{\mathrm{tail}}\|_2
\longrightarrow 0,
\]
again contradicting
\[
\|a^{(n)}-\sigma b^{(n)}\|_2^2 = 2.
\]
\end{proof}

\bibliographystyle{plain}
\bibliography{refs.bib}

\end{document}
